\documentclass[aspectratio=169]{beamer}

\usetheme{Boadilla}
\usefonttheme{professionalfonts}
\setbeamertemplate{navigation symbols}{}
\setbeamertemplate{itemize items}[circle]
\setbeamertemplate{enumerate items}[default]
\setbeamertemplate{headline}{}

\usepackage[utf8]{inputenc}
\usepackage[T1]{fontenc}
\usepackage{lmodern}
\usepackage{amsmath}
\usepackage{booktabs}
\usepackage{tabularx}
\usepackage{array}
\usepackage{hyperref}
\usepackage[strings]{underscore}
\usepackage{listings}
\usepackage{tikz}
\usetikzlibrary{arrows.meta,positioning,fit,calc,shapes.geometric}
\usepackage{xcolor}

\definecolor{courseblue}{HTML}{12355B}
\definecolor{coursegold}{HTML}{C98A2E}
\definecolor{courseice}{HTML}{EEF3F8}
\definecolor{coursegreen}{HTML}{2F6B4F}
\definecolor{coursegray}{HTML}{4B5563}
\definecolor{codebg}{HTML}{F7F9FC}
\definecolor{codecomment}{HTML}{5E6C84}
\definecolor{codestring}{HTML}{0B6E4F}

\setbeamercolor{palette primary}{bg=courseblue,fg=white}
\setbeamercolor{palette secondary}{bg=coursegold,fg=white}
\setbeamercolor{palette tertiary}{bg=courseblue,fg=white}
\setbeamercolor{palette quaternary}{bg=coursegold,fg=white}
\setbeamercolor{structure}{fg=courseblue}
\setbeamercolor{frametitle}{bg=courseice,fg=courseblue}
\setbeamercolor{title}{fg=courseblue}
\setbeamercolor{block title}{bg=courseblue,fg=white}
\setbeamercolor{block body}{bg=courseice,fg=black}
\setbeamercolor{normal text}{fg=black,bg=white}
\setbeamercolor{alerted text}{fg=coursegold}

\setbeamerfont{footline}{size=\scriptsize}

\setbeamertemplate{footline}{
  \leavevmode
  \hbox{
    \begin{beamercolorbox}[wd=.33\paperwidth,ht=2.8ex,dp=1.2ex,leftskip=0.8em]{palette primary}%
      University of Trieste
    \end{beamercolorbox}%
    \begin{beamercolorbox}[wd=.49\paperwidth,ht=2.8ex,dp=1.2ex,center]{palette tertiary}%
      \coursename{} \textbar{} \coursecode
    \end{beamercolorbox}%
    \begin{beamercolorbox}[wd=.18\paperwidth,ht=2.8ex,dp=1.2ex,rightskip=0.8em plus1fil]{palette secondary}%
      \insertframenumber{} / \inserttotalframenumber
    \end{beamercolorbox}%
  }
}

\hypersetup{
  colorlinks=true,
  linkcolor=courseblue,
  urlcolor=coursegreen
}

\lstdefinestyle{coursepython}{
  language=Python,
  basicstyle=\ttfamily\footnotesize,
  keywordstyle=\color{courseblue}\bfseries,
  commentstyle=\color{codecomment}\itshape,
  stringstyle=\color{codestring},
  showstringspaces=false,
  columns=fullflexible,
  keepspaces=true,
  frame=single,
  framerule=0.6pt,
  rulecolor=\color{courseblue},
  backgroundcolor=\color{codebg},
  xleftmargin=0.4em,
  xrightmargin=0.4em,
  aboveskip=0.4em,
  belowskip=0.2em
}

\newcommand{\coursename}{Data-driven Systems Engineering}
\newcommand{\coursecode}{440MI / 305SM}
\newcommand{\coursefooter}{}

\newcommand{\learninggoals}[1]{
  \begin{block}{Learning goals}
  #1
  \end{block}
}

\newcommand{\repoactivity}[1]{
  \begin{block}{Repo activity}
  #1
  \end{block}
}

\newcommand{\readinglist}[1]{
  \begin{block}{Papers and materials}
  {\footnotesize #1}
  \end{block}
}

\newcommand{\coursecontextframe}[5]{
  \begin{frame}{Course Map and Running Cases}
  \centering
  \begin{tikzpicture}[node distance=0.28cm]
  \node[flowbox, minimum width=1.45cm, minimum height=0.8cm] (problem) {\scriptsize Problem};
  \node[flowbox, minimum width=1.45cm, minimum height=0.8cm, right=of problem] (data) {\scriptsize Data};
  \node[flowbox, minimum width=1.45cm, minimum height=0.8cm, right=of data] (features) {\scriptsize Features};
  \node[flowbox, minimum width=1.45cm, minimum height=0.8cm, right=of features] (model) {\scriptsize Model};
  \node[flowbox, minimum width=1.45cm, minimum height=0.8cm, right=of model] (interface) {\scriptsize Interface};
  \node[flowbox, minimum width=1.45cm, minimum height=0.8cm, right=of interface] (delivery) {\scriptsize Delivery};
  \node[flowbox, minimum width=1.45cm, minimum height=0.8cm, right=of delivery] (monitor) {\scriptsize Monitor};
  \node[accentbox, minimum width=1.6cm, minimum height=0.8cm, right=of monitor] (evolution) {\scriptsize Evolution};
  \foreach \a/\b in {problem/data,data/features,features/model,model/interface,interface/delivery,delivery/monitor,monitor/evolution} {
    \draw[arrow] (\a) -- (\b);
  }
  \end{tikzpicture}

  \vspace{0.5em}
  \begin{block}{Today in the course}
  \textbf{Focus:} #1

  \textbf{Bridge:} #2
  \end{block}

  \begin{columns}[T]
  \column{0.48\textwidth}
  \begin{block}{Fraud case}
  {\footnotesize #3}
  \end{block}
  \column{0.48\textwidth}
  \begin{block}{Pasteurization case}
  {\footnotesize #4}
  \end{block}
  \end{columns}

  \vspace{0.3em}
  {\footnotesize \textbf{Course thread:} #5}
  \coursefooter
  \end{frame}
}

\newcommand{\classclosingframe}[4]{
  \begin{frame}{From This Class to the Next}
  \begin{block}{What stays with us}
  #1
  \end{block}

  \begin{columns}[T]
  \column{0.48\textwidth}
  \begin{block}{Project use}
  {\footnotesize #2}
  \end{block}
  \column{0.48\textwidth}
  \begin{block}{Next connection}
  {\footnotesize #3}
  \end{block}
  \end{columns}

  \vspace{0.3em}
  {\footnotesize \textbf{Course thread:} #4}
  \coursefooter
  \end{frame}
}

\tikzset{
  flowbox/.style={
    rounded corners=3pt,
    draw=courseblue,
    fill=courseice,
    very thick,
    minimum width=2.2cm,
    minimum height=0.9cm,
    align=center,
    font=\small
  },
  accentbox/.style={
    rounded corners=3pt,
    draw=coursegold,
    fill=coursegold!15,
    very thick,
    minimum width=2.2cm,
    minimum height=0.9cm,
    align=center,
    font=\small
  },
  arrow/.style={
    -{Latex[length=2.5mm,width=2mm]},
    thick,
    draw=courseblue
  }
}

% Visual layer on top of common.tex: tiles, icons, a progress-strip section divider, and a footer
% strip that defines the concepts appearing for the first time on each slide.
% Same look as the Class 20 deck. Load with:  \documentclass[aspectratio=169]{beamer}

\usetheme{Boadilla}
\usefonttheme{professionalfonts}
\setbeamertemplate{navigation symbols}{}
\setbeamertemplate{itemize items}[circle]
\setbeamertemplate{enumerate items}[default]
\setbeamertemplate{headline}{}

\usepackage[utf8]{inputenc}
\usepackage[T1]{fontenc}
\usepackage{lmodern}
\usepackage{amsmath}
\usepackage{booktabs}
\usepackage{tabularx}
\usepackage{array}
\usepackage{hyperref}
\usepackage[strings]{underscore}
\usepackage{listings}
\usepackage{tikz}
\usetikzlibrary{arrows.meta,positioning,fit,calc,shapes.geometric}
\usepackage{xcolor}

\definecolor{courseblue}{HTML}{12355B}
\definecolor{coursegold}{HTML}{C98A2E}
\definecolor{courseice}{HTML}{EEF3F8}
\definecolor{coursegreen}{HTML}{2F6B4F}
\definecolor{coursegray}{HTML}{4B5563}
\definecolor{codebg}{HTML}{F7F9FC}
\definecolor{codecomment}{HTML}{5E6C84}
\definecolor{codestring}{HTML}{0B6E4F}

\setbeamercolor{palette primary}{bg=courseblue,fg=white}
\setbeamercolor{palette secondary}{bg=coursegold,fg=white}
\setbeamercolor{palette tertiary}{bg=courseblue,fg=white}
\setbeamercolor{palette quaternary}{bg=coursegold,fg=white}
\setbeamercolor{structure}{fg=courseblue}
\setbeamercolor{frametitle}{bg=courseice,fg=courseblue}
\setbeamercolor{title}{fg=courseblue}
\setbeamercolor{block title}{bg=courseblue,fg=white}
\setbeamercolor{block body}{bg=courseice,fg=black}
\setbeamercolor{normal text}{fg=black,bg=white}
\setbeamercolor{alerted text}{fg=coursegold}

\setbeamerfont{footline}{size=\scriptsize}

\setbeamertemplate{footline}{
  \leavevmode
  \hbox{
    \begin{beamercolorbox}[wd=.33\paperwidth,ht=2.8ex,dp=1.2ex,leftskip=0.8em]{palette primary}%
      University of Trieste
    \end{beamercolorbox}%
    \begin{beamercolorbox}[wd=.49\paperwidth,ht=2.8ex,dp=1.2ex,center]{palette tertiary}%
      \coursename{} \textbar{} \coursecode
    \end{beamercolorbox}%
    \begin{beamercolorbox}[wd=.18\paperwidth,ht=2.8ex,dp=1.2ex,rightskip=0.8em plus1fil]{palette secondary}%
      \insertframenumber{} / \inserttotalframenumber
    \end{beamercolorbox}%
  }
}

\hypersetup{
  colorlinks=true,
  linkcolor=courseblue,
  urlcolor=coursegreen
}

\lstdefinestyle{coursepython}{
  language=Python,
  basicstyle=\ttfamily\footnotesize,
  keywordstyle=\color{courseblue}\bfseries,
  commentstyle=\color{codecomment}\itshape,
  stringstyle=\color{codestring},
  showstringspaces=false,
  columns=fullflexible,
  keepspaces=true,
  frame=single,
  framerule=0.6pt,
  rulecolor=\color{courseblue},
  backgroundcolor=\color{codebg},
  xleftmargin=0.4em,
  xrightmargin=0.4em,
  aboveskip=0.4em,
  belowskip=0.2em
}

\newcommand{\coursename}{Data-driven Systems Engineering}
\newcommand{\coursecode}{440MI / 305SM}
\newcommand{\coursefooter}{}

\newcommand{\learninggoals}[1]{
  \begin{block}{Learning goals}
  #1
  \end{block}
}

\newcommand{\repoactivity}[1]{
  \begin{block}{Repo activity}
  #1
  \end{block}
}

\newcommand{\readinglist}[1]{
  \begin{block}{Papers and materials}
  {\footnotesize #1}
  \end{block}
}

\newcommand{\coursecontextframe}[5]{
  \begin{frame}{Course Map and Running Cases}
  \centering
  \begin{tikzpicture}[node distance=0.28cm]
  \node[flowbox, minimum width=1.45cm, minimum height=0.8cm] (problem) {\scriptsize Problem};
  \node[flowbox, minimum width=1.45cm, minimum height=0.8cm, right=of problem] (data) {\scriptsize Data};
  \node[flowbox, minimum width=1.45cm, minimum height=0.8cm, right=of data] (features) {\scriptsize Features};
  \node[flowbox, minimum width=1.45cm, minimum height=0.8cm, right=of features] (model) {\scriptsize Model};
  \node[flowbox, minimum width=1.45cm, minimum height=0.8cm, right=of model] (interface) {\scriptsize Interface};
  \node[flowbox, minimum width=1.45cm, minimum height=0.8cm, right=of interface] (delivery) {\scriptsize Delivery};
  \node[flowbox, minimum width=1.45cm, minimum height=0.8cm, right=of delivery] (monitor) {\scriptsize Monitor};
  \node[accentbox, minimum width=1.6cm, minimum height=0.8cm, right=of monitor] (evolution) {\scriptsize Evolution};
  \foreach \a/\b in {problem/data,data/features,features/model,model/interface,interface/delivery,delivery/monitor,monitor/evolution} {
    \draw[arrow] (\a) -- (\b);
  }
  \end{tikzpicture}

  \vspace{0.5em}
  \begin{block}{Today in the course}
  \textbf{Focus:} #1

  \textbf{Bridge:} #2
  \end{block}

  \begin{columns}[T]
  \column{0.48\textwidth}
  \begin{block}{Fraud case}
  {\footnotesize #3}
  \end{block}
  \column{0.48\textwidth}
  \begin{block}{Pasteurization case}
  {\footnotesize #4}
  \end{block}
  \end{columns}

  \vspace{0.3em}
  {\footnotesize \textbf{Course thread:} #5}
  \coursefooter
  \end{frame}
}

\newcommand{\classclosingframe}[4]{
  \begin{frame}{From This Class to the Next}
  \begin{block}{What stays with us}
  #1
  \end{block}

  \begin{columns}[T]
  \column{0.48\textwidth}
  \begin{block}{Project use}
  {\footnotesize #2}
  \end{block}
  \column{0.48\textwidth}
  \begin{block}{Next connection}
  {\footnotesize #3}
  \end{block}
  \end{columns}

  \vspace{0.3em}
  {\footnotesize \textbf{Course thread:} #4}
  \coursefooter
  \end{frame}
}

\tikzset{
  flowbox/.style={
    rounded corners=3pt,
    draw=courseblue,
    fill=courseice,
    very thick,
    minimum width=2.2cm,
    minimum height=0.9cm,
    align=center,
    font=\small
  },
  accentbox/.style={
    rounded corners=3pt,
    draw=coursegold,
    fill=coursegold!15,
    very thick,
    minimum width=2.2cm,
    minimum height=0.9cm,
    align=center,
    font=\small
  },
  arrow/.style={
    -{Latex[length=2.5mm,width=2mm]},
    thick,
    draw=courseblue
  }
}
  % Visual layer on top of common.tex: tiles, icons, a progress-strip section divider, and a footer
% strip that defines the concepts appearing for the first time on each slide.
% Same look as the Class 20 deck. Load with:  \documentclass[aspectratio=169]{beamer}

\usetheme{Boadilla}
\usefonttheme{professionalfonts}
\setbeamertemplate{navigation symbols}{}
\setbeamertemplate{itemize items}[circle]
\setbeamertemplate{enumerate items}[default]
\setbeamertemplate{headline}{}

\usepackage[utf8]{inputenc}
\usepackage[T1]{fontenc}
\usepackage{lmodern}
\usepackage{amsmath}
\usepackage{booktabs}
\usepackage{tabularx}
\usepackage{array}
\usepackage{hyperref}
\usepackage[strings]{underscore}
\usepackage{listings}
\usepackage{tikz}
\usetikzlibrary{arrows.meta,positioning,fit,calc,shapes.geometric}
\usepackage{xcolor}

\definecolor{courseblue}{HTML}{12355B}
\definecolor{coursegold}{HTML}{C98A2E}
\definecolor{courseice}{HTML}{EEF3F8}
\definecolor{coursegreen}{HTML}{2F6B4F}
\definecolor{coursegray}{HTML}{4B5563}
\definecolor{codebg}{HTML}{F7F9FC}
\definecolor{codecomment}{HTML}{5E6C84}
\definecolor{codestring}{HTML}{0B6E4F}

\setbeamercolor{palette primary}{bg=courseblue,fg=white}
\setbeamercolor{palette secondary}{bg=coursegold,fg=white}
\setbeamercolor{palette tertiary}{bg=courseblue,fg=white}
\setbeamercolor{palette quaternary}{bg=coursegold,fg=white}
\setbeamercolor{structure}{fg=courseblue}
\setbeamercolor{frametitle}{bg=courseice,fg=courseblue}
\setbeamercolor{title}{fg=courseblue}
\setbeamercolor{block title}{bg=courseblue,fg=white}
\setbeamercolor{block body}{bg=courseice,fg=black}
\setbeamercolor{normal text}{fg=black,bg=white}
\setbeamercolor{alerted text}{fg=coursegold}

\setbeamerfont{footline}{size=\scriptsize}

\setbeamertemplate{footline}{
  \leavevmode
  \hbox{
    \begin{beamercolorbox}[wd=.33\paperwidth,ht=2.8ex,dp=1.2ex,leftskip=0.8em]{palette primary}%
      University of Trieste
    \end{beamercolorbox}%
    \begin{beamercolorbox}[wd=.49\paperwidth,ht=2.8ex,dp=1.2ex,center]{palette tertiary}%
      \coursename{} \textbar{} \coursecode
    \end{beamercolorbox}%
    \begin{beamercolorbox}[wd=.18\paperwidth,ht=2.8ex,dp=1.2ex,rightskip=0.8em plus1fil]{palette secondary}%
      \insertframenumber{} / \inserttotalframenumber
    \end{beamercolorbox}%
  }
}

\hypersetup{
  colorlinks=true,
  linkcolor=courseblue,
  urlcolor=coursegreen
}

\lstdefinestyle{coursepython}{
  language=Python,
  basicstyle=\ttfamily\footnotesize,
  keywordstyle=\color{courseblue}\bfseries,
  commentstyle=\color{codecomment}\itshape,
  stringstyle=\color{codestring},
  showstringspaces=false,
  columns=fullflexible,
  keepspaces=true,
  frame=single,
  framerule=0.6pt,
  rulecolor=\color{courseblue},
  backgroundcolor=\color{codebg},
  xleftmargin=0.4em,
  xrightmargin=0.4em,
  aboveskip=0.4em,
  belowskip=0.2em
}

\newcommand{\coursename}{Data-driven Systems Engineering}
\newcommand{\coursecode}{440MI / 305SM}
\newcommand{\coursefooter}{}

\newcommand{\learninggoals}[1]{
  \begin{block}{Learning goals}
  #1
  \end{block}
}

\newcommand{\repoactivity}[1]{
  \begin{block}{Repo activity}
  #1
  \end{block}
}

\newcommand{\readinglist}[1]{
  \begin{block}{Papers and materials}
  {\footnotesize #1}
  \end{block}
}

\newcommand{\coursecontextframe}[5]{
  \begin{frame}{Course Map and Running Cases}
  \centering
  \begin{tikzpicture}[node distance=0.28cm]
  \node[flowbox, minimum width=1.45cm, minimum height=0.8cm] (problem) {\scriptsize Problem};
  \node[flowbox, minimum width=1.45cm, minimum height=0.8cm, right=of problem] (data) {\scriptsize Data};
  \node[flowbox, minimum width=1.45cm, minimum height=0.8cm, right=of data] (features) {\scriptsize Features};
  \node[flowbox, minimum width=1.45cm, minimum height=0.8cm, right=of features] (model) {\scriptsize Model};
  \node[flowbox, minimum width=1.45cm, minimum height=0.8cm, right=of model] (interface) {\scriptsize Interface};
  \node[flowbox, minimum width=1.45cm, minimum height=0.8cm, right=of interface] (delivery) {\scriptsize Delivery};
  \node[flowbox, minimum width=1.45cm, minimum height=0.8cm, right=of delivery] (monitor) {\scriptsize Monitor};
  \node[accentbox, minimum width=1.6cm, minimum height=0.8cm, right=of monitor] (evolution) {\scriptsize Evolution};
  \foreach \a/\b in {problem/data,data/features,features/model,model/interface,interface/delivery,delivery/monitor,monitor/evolution} {
    \draw[arrow] (\a) -- (\b);
  }
  \end{tikzpicture}

  \vspace{0.5em}
  \begin{block}{Today in the course}
  \textbf{Focus:} #1

  \textbf{Bridge:} #2
  \end{block}

  \begin{columns}[T]
  \column{0.48\textwidth}
  \begin{block}{Fraud case}
  {\footnotesize #3}
  \end{block}
  \column{0.48\textwidth}
  \begin{block}{Pasteurization case}
  {\footnotesize #4}
  \end{block}
  \end{columns}

  \vspace{0.3em}
  {\footnotesize \textbf{Course thread:} #5}
  \coursefooter
  \end{frame}
}

\newcommand{\classclosingframe}[4]{
  \begin{frame}{From This Class to the Next}
  \begin{block}{What stays with us}
  #1
  \end{block}

  \begin{columns}[T]
  \column{0.48\textwidth}
  \begin{block}{Project use}
  {\footnotesize #2}
  \end{block}
  \column{0.48\textwidth}
  \begin{block}{Next connection}
  {\footnotesize #3}
  \end{block}
  \end{columns}

  \vspace{0.3em}
  {\footnotesize \textbf{Course thread:} #4}
  \coursefooter
  \end{frame}
}

\tikzset{
  flowbox/.style={
    rounded corners=3pt,
    draw=courseblue,
    fill=courseice,
    very thick,
    minimum width=2.2cm,
    minimum height=0.9cm,
    align=center,
    font=\small
  },
  accentbox/.style={
    rounded corners=3pt,
    draw=coursegold,
    fill=coursegold!15,
    very thick,
    minimum width=2.2cm,
    minimum height=0.9cm,
    align=center,
    font=\small
  },
  arrow/.style={
    -{Latex[length=2.5mm,width=2mm]},
    thick,
    draw=courseblue
  }
}
  % Visual layer on top of common.tex: tiles, icons, a progress-strip section divider, and a footer
% strip that defines the concepts appearing for the first time on each slide.
% Same look as the Class 20 deck. Load with:  \input{common}  \input{common_visual}
%
% Before the first \sectiondivider, a deck must define
%   \sectioncount   the number of sections, e.g. 4
%   \sectionlabels  a \foreach list of number/label pairs, e.g. 1/Need,2/Testable,3/Manage,4/Practice
% Icons come from the fontawesome5 package (Font Awesome 5 Free), e.g. \faIcon{robot}.

\usepackage{colortbl}
\usepackage{fontawesome5}
\usetikzlibrary{backgrounds}

% Slides read better without automatic hyphenation.
\hyphenpenalty=10000
\exhyphenpenalty=10000

\definecolor{coursered}{HTML}{A23B2A}
\newcommand{\goldink}{\color{coursegold!65!black}}

\tikzset{
  tile/.style={rounded corners=4pt, fill=courseice, draw=none, align=center, font=\footnotesize, inner sep=5pt},
  tileblue/.style={tile, fill=white, draw=courseblue!55, thick},
  tilegold/.style={tile, fill=coursegold!14, draw=coursegold, thick},
  tilegreen/.style={tile, fill=coursegreen!10, draw=coursegreen!70, thick},
  tilered/.style={tile, fill=coursered!8, draw=coursered!70, thick},
  tilegray/.style={tile, fill=white, draw=coursegray!55, thick, dashed},
  tag/.style={rounded corners=2pt, fill=coursegold, text=white, font=\bfseries\tiny, inner sep=2.5pt},
  chip/.style={rounded corners=3pt, fill=courseblue, text=white, font=\footnotesize\bfseries, align=center, inner sep=4pt},
  note/.style={font=\scriptsize, text=coursegray, align=left},
  arrowgold/.style={-{Latex[length=2.5mm,width=2mm]}, thick, draw=coursegold},
  arrowred/.style={-{Latex[length=2.5mm,width=2mm]}, thick, draw=coursered},
  arrowgreen/.style={-{Latex[length=2.5mm,width=2mm]}, thick, draw=coursegreen},
  arrowgray/.style={-{Latex[length=2mm,width=1.6mm]}, thick, draw=coursegray!70},
}

% A round icon badge, drawn inside a tikzpicture: x y colour icon
\newcommand{\badge}[4]{%
  \node[circle, fill=#3, text=white, minimum size=0.85cm, inner sep=0pt] at (#1,#2) {\large\faIcon{#4}};
}

% Inline tags for the kind of requirement.
\newcommand{\Ftag}{\colorbox{courseblue}{\color{white}\bfseries\scriptsize\ \ F\ \ }}
\newcommand{\NFtag}{\colorbox{coursegold}{\color{white}\bfseries\scriptsize\ NF\ }}
\newcommand{\Dtag}{\colorbox{coursegreen}{\color{white}\bfseries\scriptsize\ \ D\ \ }}
\newcommand{\MLtag}{\colorbox{coursered}{\color{white}\bfseries\scriptsize\ ML\ }}

% Footer: a strip with the concepts that appear for the first time on the slide, above the
% usual bar. The bar total counts only the main lecture; a backup slide is numbered separately.
\makeatletter
\def\conceptstext{}
\def\currentsectionname{}
\setbeamercolor{conceptbar}{bg=coursegold!12,fg=black}
\newcommand{\cterm}[2]{{\bfseries\color{courseblue}#1:}\ #2\hspace{0.8em}}
\newcommand{\newtag}{\colorbox{coursegold}{\color{white}\bfseries\tiny\ NEW\ }\hspace{0.6em}}
\setbeamertemplate{footline}{
  \leavevmode
  \vbox{%
    \ifx\conceptstext\empty\else
      \hbox{\color{coursegold}\rule{\paperwidth}{0.8pt}}%
      \hbox{\begin{beamercolorbox}[wd=\paperwidth,leftskip=0.5cm,rightskip=0.5cm,sep=3pt]{conceptbar}\scriptsize\newtag\conceptstext\end{beamercolorbox}}%
    \fi
    \hbox{%
      \begin{beamercolorbox}[wd=.33\paperwidth,ht=2.8ex,dp=1.2ex,leftskip=0.8em]{palette primary}%
        \ifbeamer@inappendix Backup\else\ifnum\value{section}>0 \arabic{section}\ \textbar\ \currentsectionname\else University of Trieste\fi\fi
      \end{beamercolorbox}%
      \begin{beamercolorbox}[wd=.49\paperwidth,ht=2.8ex,dp=1.2ex,center]{palette tertiary}%
        \coursename{} \textbar{} \coursecode
      \end{beamercolorbox}%
      \begin{beamercolorbox}[wd=.18\paperwidth,ht=2.8ex,dp=1.2ex,right,rightskip=0.8em]{palette secondary}%
        \ifbeamer@inappendix Backup \insertframenumberinappendix\else\insertframenumber{} / \insertmainframenumber\fi
      \end{beamercolorbox}%
    }}}
\makeatother
\newenvironment{cframe}[2]{\gdef\conceptstext{#2}\begin{frame}{#1}}{\end{frame}\gdef\conceptstext{}}

% Section divider: progress strip, recap of the previous part, and the roadmap of the next one.
% Arguments: number, title, icon, previously, items.
\newcommand{\sectiondivider}[5]{%
  \section{#2}\gdef\currentsectionname{#2}
  \begin{frame}
  \centering
  \begin{tikzpicture}
  \pgfmathsetmacro{\chipw}{13.6/\sectioncount}
  \foreach \i/\lab in \sectionlabels {
    \ifnum\i<#1 \def\chipfill{courseblue!22}\def\chiptext{courseblue}\else
    \ifnum\i=#1 \def\chipfill{courseblue}\def\chiptext{white}\else
    \def\chipfill{courseice}\def\chiptext{coursegray}\fi\fi
    \node[rounded corners=3pt, minimum width={(\chipw-0.1)*1cm}, minimum height=0.6cm, font=\scriptsize\bfseries, fill=\chipfill, text=\chiptext] at ({\chipw/2+(\i-1)*\chipw},0) {\i\ \lab};
  }
  \node[circle, fill=courseblue, text=white, font=\Huge\bfseries, minimum size=1.7cm, inner sep=0] at (1.0,-1.7) {#1};
  \node[anchor=west, font=\huge\bfseries, text=courseblue, text width=8.8cm, align=left] at (2.2,-1.7) {#2};
  \node[circle, fill=coursegold!16, draw=coursegold, thick, text=coursegold!75!black, font=\fontsize{26}{26}\selectfont, minimum size=1.9cm, inner sep=0] at (12.6,-1.7) {\faIcon{#3}};
  \ifx\relax#4\relax\else
    \node[anchor=north west, font=\small\itshape, text=coursegray, text width=13.2cm, align=left] at (0.1,-2.95) {Previously: #4};
  \fi
  \node[anchor=north west, font=\footnotesize\bfseries, text=coursegold!65!black] at (0.1,-3.95) {IN THIS PART};
  \node[anchor=north west, font=\small, text width=13.2cm, align=left] at (0.1,-4.35) {#5};
  \end{tikzpicture}
  \end{frame}
}

%
% Before the first \sectiondivider, a deck must define
%   \sectioncount   the number of sections, e.g. 4
%   \sectionlabels  a \foreach list of number/label pairs, e.g. 1/Need,2/Testable,3/Manage,4/Practice
% Icons come from the fontawesome5 package (Font Awesome 5 Free), e.g. \faIcon{robot}.

\usepackage{colortbl}
\usepackage{fontawesome5}
\usetikzlibrary{backgrounds}

% Slides read better without automatic hyphenation.
\hyphenpenalty=10000
\exhyphenpenalty=10000

\definecolor{coursered}{HTML}{A23B2A}
\newcommand{\goldink}{\color{coursegold!65!black}}

\tikzset{
  tile/.style={rounded corners=4pt, fill=courseice, draw=none, align=center, font=\footnotesize, inner sep=5pt},
  tileblue/.style={tile, fill=white, draw=courseblue!55, thick},
  tilegold/.style={tile, fill=coursegold!14, draw=coursegold, thick},
  tilegreen/.style={tile, fill=coursegreen!10, draw=coursegreen!70, thick},
  tilered/.style={tile, fill=coursered!8, draw=coursered!70, thick},
  tilegray/.style={tile, fill=white, draw=coursegray!55, thick, dashed},
  tag/.style={rounded corners=2pt, fill=coursegold, text=white, font=\bfseries\tiny, inner sep=2.5pt},
  chip/.style={rounded corners=3pt, fill=courseblue, text=white, font=\footnotesize\bfseries, align=center, inner sep=4pt},
  note/.style={font=\scriptsize, text=coursegray, align=left},
  arrowgold/.style={-{Latex[length=2.5mm,width=2mm]}, thick, draw=coursegold},
  arrowred/.style={-{Latex[length=2.5mm,width=2mm]}, thick, draw=coursered},
  arrowgreen/.style={-{Latex[length=2.5mm,width=2mm]}, thick, draw=coursegreen},
  arrowgray/.style={-{Latex[length=2mm,width=1.6mm]}, thick, draw=coursegray!70},
}

% A round icon badge, drawn inside a tikzpicture: x y colour icon
\newcommand{\badge}[4]{%
  \node[circle, fill=#3, text=white, minimum size=0.85cm, inner sep=0pt] at (#1,#2) {\large\faIcon{#4}};
}

% Inline tags for the kind of requirement.
\newcommand{\Ftag}{\colorbox{courseblue}{\color{white}\bfseries\scriptsize\ \ F\ \ }}
\newcommand{\NFtag}{\colorbox{coursegold}{\color{white}\bfseries\scriptsize\ NF\ }}
\newcommand{\Dtag}{\colorbox{coursegreen}{\color{white}\bfseries\scriptsize\ \ D\ \ }}
\newcommand{\MLtag}{\colorbox{coursered}{\color{white}\bfseries\scriptsize\ ML\ }}

% Footer: a strip with the concepts that appear for the first time on the slide, above the
% usual bar. The bar total counts only the main lecture; a backup slide is numbered separately.
\makeatletter
\def\conceptstext{}
\def\currentsectionname{}
\setbeamercolor{conceptbar}{bg=coursegold!12,fg=black}
\newcommand{\cterm}[2]{{\bfseries\color{courseblue}#1:}\ #2\hspace{0.8em}}
\newcommand{\newtag}{\colorbox{coursegold}{\color{white}\bfseries\tiny\ NEW\ }\hspace{0.6em}}
\setbeamertemplate{footline}{
  \leavevmode
  \vbox{%
    \ifx\conceptstext\empty\else
      \hbox{\color{coursegold}\rule{\paperwidth}{0.8pt}}%
      \hbox{\begin{beamercolorbox}[wd=\paperwidth,leftskip=0.5cm,rightskip=0.5cm,sep=3pt]{conceptbar}\scriptsize\newtag\conceptstext\end{beamercolorbox}}%
    \fi
    \hbox{%
      \begin{beamercolorbox}[wd=.33\paperwidth,ht=2.8ex,dp=1.2ex,leftskip=0.8em]{palette primary}%
        \ifbeamer@inappendix Backup\else\ifnum\value{section}>0 \arabic{section}\ \textbar\ \currentsectionname\else University of Trieste\fi\fi
      \end{beamercolorbox}%
      \begin{beamercolorbox}[wd=.49\paperwidth,ht=2.8ex,dp=1.2ex,center]{palette tertiary}%
        \coursename{} \textbar{} \coursecode
      \end{beamercolorbox}%
      \begin{beamercolorbox}[wd=.18\paperwidth,ht=2.8ex,dp=1.2ex,right,rightskip=0.8em]{palette secondary}%
        \ifbeamer@inappendix Backup \insertframenumberinappendix\else\insertframenumber{} / \insertmainframenumber\fi
      \end{beamercolorbox}%
    }}}
\makeatother
\newenvironment{cframe}[2]{\gdef\conceptstext{#2}\begin{frame}{#1}}{\end{frame}\gdef\conceptstext{}}

% Section divider: progress strip, recap of the previous part, and the roadmap of the next one.
% Arguments: number, title, icon, previously, items.
\newcommand{\sectiondivider}[5]{%
  \section{#2}\gdef\currentsectionname{#2}
  \begin{frame}
  \centering
  \begin{tikzpicture}
  \pgfmathsetmacro{\chipw}{13.6/\sectioncount}
  \foreach \i/\lab in \sectionlabels {
    \ifnum\i<#1 \def\chipfill{courseblue!22}\def\chiptext{courseblue}\else
    \ifnum\i=#1 \def\chipfill{courseblue}\def\chiptext{white}\else
    \def\chipfill{courseice}\def\chiptext{coursegray}\fi\fi
    \node[rounded corners=3pt, minimum width={(\chipw-0.1)*1cm}, minimum height=0.6cm, font=\scriptsize\bfseries, fill=\chipfill, text=\chiptext] at ({\chipw/2+(\i-1)*\chipw},0) {\i\ \lab};
  }
  \node[circle, fill=courseblue, text=white, font=\Huge\bfseries, minimum size=1.7cm, inner sep=0] at (1.0,-1.7) {#1};
  \node[anchor=west, font=\huge\bfseries, text=courseblue, text width=8.8cm, align=left] at (2.2,-1.7) {#2};
  \node[circle, fill=coursegold!16, draw=coursegold, thick, text=coursegold!75!black, font=\fontsize{26}{26}\selectfont, minimum size=1.9cm, inner sep=0] at (12.6,-1.7) {\faIcon{#3}};
  \ifx\relax#4\relax\else
    \node[anchor=north west, font=\small\itshape, text=coursegray, text width=13.2cm, align=left] at (0.1,-2.95) {Previously: #4};
  \fi
  \node[anchor=north west, font=\footnotesize\bfseries, text=coursegold!65!black] at (0.1,-3.95) {IN THIS PART};
  \node[anchor=north west, font=\small, text width=13.2cm, align=left] at (0.1,-4.35) {#5};
  \end{tikzpicture}
  \end{frame}
}

%
% Before the first \sectiondivider, a deck must define
%   \sectioncount   the number of sections, e.g. 4
%   \sectionlabels  a \foreach list of number/label pairs, e.g. 1/Need,2/Testable,3/Manage,4/Practice
% Icons come from the fontawesome5 package (Font Awesome 5 Free), e.g. \faIcon{robot}.

\usepackage{colortbl}
\usepackage{fontawesome5}
\usetikzlibrary{backgrounds}

% Slides read better without automatic hyphenation.
\hyphenpenalty=10000
\exhyphenpenalty=10000

\definecolor{coursered}{HTML}{A23B2A}
\newcommand{\goldink}{\color{coursegold!65!black}}

\tikzset{
  tile/.style={rounded corners=4pt, fill=courseice, draw=none, align=center, font=\footnotesize, inner sep=5pt},
  tileblue/.style={tile, fill=white, draw=courseblue!55, thick},
  tilegold/.style={tile, fill=coursegold!14, draw=coursegold, thick},
  tilegreen/.style={tile, fill=coursegreen!10, draw=coursegreen!70, thick},
  tilered/.style={tile, fill=coursered!8, draw=coursered!70, thick},
  tilegray/.style={tile, fill=white, draw=coursegray!55, thick, dashed},
  tag/.style={rounded corners=2pt, fill=coursegold, text=white, font=\bfseries\tiny, inner sep=2.5pt},
  chip/.style={rounded corners=3pt, fill=courseblue, text=white, font=\footnotesize\bfseries, align=center, inner sep=4pt},
  note/.style={font=\scriptsize, text=coursegray, align=left},
  arrowgold/.style={-{Latex[length=2.5mm,width=2mm]}, thick, draw=coursegold},
  arrowred/.style={-{Latex[length=2.5mm,width=2mm]}, thick, draw=coursered},
  arrowgreen/.style={-{Latex[length=2.5mm,width=2mm]}, thick, draw=coursegreen},
  arrowgray/.style={-{Latex[length=2mm,width=1.6mm]}, thick, draw=coursegray!70},
}

% A round icon badge, drawn inside a tikzpicture: x y colour icon
\newcommand{\badge}[4]{%
  \node[circle, fill=#3, text=white, minimum size=0.85cm, inner sep=0pt] at (#1,#2) {\large\faIcon{#4}};
}

% Inline tags for the kind of requirement.
\newcommand{\Ftag}{\colorbox{courseblue}{\color{white}\bfseries\scriptsize\ \ F\ \ }}
\newcommand{\NFtag}{\colorbox{coursegold}{\color{white}\bfseries\scriptsize\ NF\ }}
\newcommand{\Dtag}{\colorbox{coursegreen}{\color{white}\bfseries\scriptsize\ \ D\ \ }}
\newcommand{\MLtag}{\colorbox{coursered}{\color{white}\bfseries\scriptsize\ ML\ }}

% Footer: a strip with the concepts that appear for the first time on the slide, above the
% usual bar. The bar total counts only the main lecture; a backup slide is numbered separately.
\makeatletter
\def\conceptstext{}
\def\currentsectionname{}
\setbeamercolor{conceptbar}{bg=coursegold!12,fg=black}
\newcommand{\cterm}[2]{{\bfseries\color{courseblue}#1:}\ #2\hspace{0.8em}}
\newcommand{\newtag}{\colorbox{coursegold}{\color{white}\bfseries\tiny\ NEW\ }\hspace{0.6em}}
\setbeamertemplate{footline}{
  \leavevmode
  \vbox{%
    \ifx\conceptstext\empty\else
      \hbox{\color{coursegold}\rule{\paperwidth}{0.8pt}}%
      \hbox{\begin{beamercolorbox}[wd=\paperwidth,leftskip=0.5cm,rightskip=0.5cm,sep=3pt]{conceptbar}\scriptsize\newtag\conceptstext\end{beamercolorbox}}%
    \fi
    \hbox{%
      \begin{beamercolorbox}[wd=.33\paperwidth,ht=2.8ex,dp=1.2ex,leftskip=0.8em]{palette primary}%
        \ifbeamer@inappendix Backup\else\ifnum\value{section}>0 \arabic{section}\ \textbar\ \currentsectionname\else University of Trieste\fi\fi
      \end{beamercolorbox}%
      \begin{beamercolorbox}[wd=.49\paperwidth,ht=2.8ex,dp=1.2ex,center]{palette tertiary}%
        \coursename{} \textbar{} \coursecode
      \end{beamercolorbox}%
      \begin{beamercolorbox}[wd=.18\paperwidth,ht=2.8ex,dp=1.2ex,right,rightskip=0.8em]{palette secondary}%
        \ifbeamer@inappendix Backup \insertframenumberinappendix\else\insertframenumber{} / \insertmainframenumber\fi
      \end{beamercolorbox}%
    }}}
\makeatother
\newenvironment{cframe}[2]{\gdef\conceptstext{#2}\begin{frame}{#1}}{\end{frame}\gdef\conceptstext{}}

% Section divider: progress strip, recap of the previous part, and the roadmap of the next one.
% Arguments: number, title, icon, previously, items.
\newcommand{\sectiondivider}[5]{%
  \section{#2}\gdef\currentsectionname{#2}
  \begin{frame}
  \centering
  \begin{tikzpicture}
  \pgfmathsetmacro{\chipw}{13.6/\sectioncount}
  \foreach \i/\lab in \sectionlabels {
    \ifnum\i<#1 \def\chipfill{courseblue!22}\def\chiptext{courseblue}\else
    \ifnum\i=#1 \def\chipfill{courseblue}\def\chiptext{white}\else
    \def\chipfill{courseice}\def\chiptext{coursegray}\fi\fi
    \node[rounded corners=3pt, minimum width={(\chipw-0.1)*1cm}, minimum height=0.6cm, font=\scriptsize\bfseries, fill=\chipfill, text=\chiptext] at ({\chipw/2+(\i-1)*\chipw},0) {\i\ \lab};
  }
  \node[circle, fill=courseblue, text=white, font=\Huge\bfseries, minimum size=1.7cm, inner sep=0] at (1.0,-1.7) {#1};
  \node[anchor=west, font=\huge\bfseries, text=courseblue, text width=8.8cm, align=left] at (2.2,-1.7) {#2};
  \node[circle, fill=coursegold!16, draw=coursegold, thick, text=coursegold!75!black, font=\fontsize{26}{26}\selectfont, minimum size=1.9cm, inner sep=0] at (12.6,-1.7) {\faIcon{#3}};
  \ifx\relax#4\relax\else
    \node[anchor=north west, font=\small\itshape, text=coursegray, text width=13.2cm, align=left] at (0.1,-2.95) {Previously: #4};
  \fi
  \node[anchor=north west, font=\footnotesize\bfseries, text=coursegold!65!black] at (0.1,-3.95) {IN THIS PART};
  \node[anchor=north west, font=\small, text width=13.2cm, align=left] at (0.1,-4.35) {#5};
  \end{tikzpicture}
  \end{frame}
}

\usepackage{amssymb}
\usetikzlibrary{decorations.pathreplacing}

% Four parts; the progress strip on each divider uses these labels.
\newcommand{\sectioncount}{4}
\newcommand{\sectionlabels}{1/Online learning,2/Concept drift,3/Detection,4/Adaptation}

\title{Class 29: Online Learning and Concept Drift}
\author{\coursename}
\date{}

% ------------------------------------------------------------------ helpers
% A gold strip with a labelled idea: y icon label text
\newcommand{\callout}[4]{%
  \node[tilegold, anchor=north west, minimum width=13.4cm, minimum height=0.8cm, text width=12.6cm, align=left, font=\footnotesize] at (0,#1) {{\bfseries\goldink\faIcon{#2}\ #3}\ #4};
}
\newcommand{\keyidea}[2]{\callout{#1}{lightbulb}{Key idea.}{#2}}
% A red strip: y icon label text
\newcommand{\warning}[4]{%
  \node[tilered, anchor=north west, minimum width=13.4cm, minimum height=0.8cm, text width=12.6cm, align=left, font=\footnotesize] at (0,#1) {{\bfseries\color{coursered}\faIcon{#2}\ #3}\ #4};
}
% An information tile: x y width height style colour icon title text
\newcommand{\infotile}[9]{%
  \node[#5, anchor=north west, minimum width=#3cm, minimum height=#4cm] at (#1,#2) {};
  \badge{#1+0.55}{#2-0.55}{#6}{#7}
  \node[anchor=west, text width={(#3-1.35)*1cm}, align=left, font=\footnotesize\bfseries, text=courseblue] at (#1+1.05,#2-0.55) {#8};
  \node[anchor=north west, text width={(#3-0.4)*1cm}, align=left, font=\scriptsize, inner sep=0pt] at (#1+0.2,#2-1.0) {#9};
}
% One row of the repository map: y path text stage-colour stage
\newcommand{\reporow}[5]{%
  \node[tile, fill=courseice, anchor=north west, minimum width=4.6cm, minimum height=0.74cm, inner sep=2pt, text width=4.3cm, align=left, font=\ttfamily\scriptsize] at (0,#1) {#2};
  \node[tile, fill=white, draw=courseblue!40, thick, anchor=north west, minimum width=6.3cm, minimum height=0.74cm, inner sep=2pt, text width=5.9cm, align=left, font=\scriptsize] at (4.75,#1) {#3};
  \node[chip, fill=#4, anchor=north west, minimum width=2.15cm, minimum height=0.74cm, font=\scriptsize\bfseries] at (11.25,#1) {#5};
}
% A reference row: y height icon label text
\newcommand{\refrow}[5]{%
  \node[chip, anchor=west, minimum width=3.4cm, minimum height=0.68cm, font=\scriptsize\bfseries, align=left] at (0,#1-#2/2) {\faIcon{#3}\ #4};
  \node[anchor=west, text width=9.8cm, font=\scriptsize, align=left] at (3.6,#1-#2/2) {#5};
  \draw[courseice!70!courseblue, line width=0.5pt] (0,#1-#2) -- (13.7,#1-#2);
}
% A row of a two-alternative comparison: y icon criterion text-left text-right
\newcommand{\comparerow}[5]{%
  \node[chip, fill=courseice, text=courseblue, anchor=north west, minimum width=3.0cm, minimum height=0.7cm, font=\scriptsize\bfseries, align=left] at (0,#1) {\faIcon{#2}\ #3};
  \node[tile, fill=white, draw=courseblue!60, thick, anchor=north west, minimum width=5.05cm, minimum height=0.7cm, text width=4.8cm, align=left, font=\scriptsize, inner sep=2pt] at (3.15,#1) {#4};
  \node[tile, fill=white, draw=coursegreen!70, thick, anchor=north west, minimum width=5.05cm, minimum height=0.7cm, text width=4.8cm, align=left, font=\scriptsize, inner sep=2pt] at (8.35,#1) {#5};
}

\begin{document}

% ------------------------------------------------------------------ title
\frame{\titlepage \coursefooter}

% ------------------------------------------------------------------ roadmap
\begin{cframe}{Roadmap and goals}{\cterm{Online learning}{updating a model one observation at a time, as the data arrive.}\cterm{Concept drift}{a change over time of the data-generating process.}}
\centering
\begin{tikzpicture}
\foreach \i/\col/\ico/\ttl/\txt in {0/courseblue/stream/Online learning/{Batch versus incremental models; prequential evaluation.}, 1/courseblue/random/Concept drift/{Virtual and real drift; four temporal patterns.}, 2/coursered/search/Detection/{ADWIN on the error stream, and drift reports.}, 3/coursegreen/sync-alt/Adaptation/{What the model does once drift is detected.}} {
  \pgfmathsetmacro{\x}{\i*3.4}
  \node[tile, fill=white, draw=\col, thick, anchor=north west, minimum width=3.2cm, minimum height=2.55cm] at (\x,0) {};
  \node[circle, fill=\col, text=white, minimum size=0.5cm, inner sep=0, font=\small\bfseries] at (\x+0.4,-0.45) {\the\numexpr\i+1\relax};
  \node[anchor=west, text width=1.7cm, align=left, font=\footnotesize\bfseries, text=courseblue, inner sep=0pt] at (\x+0.8,-0.45) {\ttl};
  \node[text=\col!70, font=\Large] at (\x+2.8,-0.45) {\faIcon{\ico}};
  \node[anchor=north west, text width=2.8cm, align=left, font=\scriptsize, inner sep=0pt] at (\x+0.2,-1.0) {\txt};
}
\node[anchor=west, font=\tiny\bfseries, text=coursegold!65!black] at (0,-2.85) {WHAT YOU WILL BE ABLE TO DO};
\foreach \x/\y/\ico/\t in {0/-3.0/balance-scale/{Choose between batch retraining and online learning for a given system.}, 6.8/-3.0/redo-alt/{Evaluate a streaming model prequentially, without a held-out set.}, 0/-3.75/random/{Distinguish virtual from real drift, and the four temporal patterns.}, 6.8/-3.75/search/{Explain how ADWIN decides that the mean of a stream has changed.}, 0/-4.5/chart-line/{Read a detection delay and a false-alarm count from an experiment.}, 6.8/-4.5/sync-alt/{Select an adaptation strategy, and avoid the reset loop.}} {
  \node[circle, fill=courseblue, text=white, minimum size=0.45cm, inner sep=0, font=\scriptsize] at (\x+0.3,\y-0.3) {\faIcon{\ico}};
  \node[anchor=west, text width=5.8cm, align=left, font=\scriptsize, inner sep=0pt] at (\x+0.7,\y-0.3) {\t};
}
\end{tikzpicture}
\coursefooter
\end{cframe}

% ------------------------------------------------------------------ course context
\begin{cframe}{Where this class sits in the course}{\cterm{Running case}{one of two examples, fraud scoring and pasteurization monitoring, that return in every class.}}
\centering
\begin{tikzpicture}
\foreach \i/\lab/\bg/\fg in {0/Problem/courseice/coursegray, 1/Data/courseice/coursegray, 2/Features/courseice/coursegray, 3/Model/courseblue/white, 4/Interface/courseice/coursegray, 5/Delivery/courseice/coursegray, 6/Monitor/courseblue/white, 7/Evolution/courseblue/white} {
  \pgfmathsetmacro{\x}{\i*1.68}
  \node[rounded corners=3pt, fill=\bg, text=\fg, anchor=north west, minimum width=1.55cm, minimum height=0.6cm, font=\scriptsize\bfseries] at (\x,0) {\lab};
}
\node[tilegold, anchor=north west, minimum width=6.6cm, minimum height=1.3cm, text width=6.2cm, align=left] at (0,-0.9) {{\bfseries\goldink\faIcon{bullseye}\ Focus.} Models that keep learning after deployment, and the statistical detection of change.};
\node[tileblue, anchor=north west, minimum width=6.6cm, minimum height=1.3cm, text width=6.2cm, align=left] at (6.8,-0.9) {{\bfseries\color{courseblue}\faIcon{link}\ Bridge.} Class 27 named drift as a trigger for retraining; this class opens the trigger, and Class 30 builds it into the challenge.};
\node[tileblue, anchor=north west, minimum width=6.6cm, minimum height=1.5cm, text width=6.2cm, align=left, font=\scriptsize] at (0,-2.5) {{\bfseries\color{courseblue}\faIcon{credit-card}\ Fraud case.} Seasonal spending and new merchant categories shift the inputs; new fraud tactics change the relation between inputs and label.};
\node[tileblue, anchor=north west, minimum width=6.6cm, minimum height=1.5cm, text width=6.2cm, align=left, font=\scriptsize] at (6.8,-2.5) {{\bfseries\color{courseblue}\faIcon{industry}\ Pasteurization case.} Sensor recalibration is sudden; raw-material variability and equipment wear are gradual or incremental.};
\keyidea{-4.2}{A validated model is valid for the process that generated its data. When the process changes, validity must be re-established, preferably as soon as the change occurs.}
\end{tikzpicture}
\coursefooter
\end{cframe}

% ================================================================== PART 1
\sectiondivider{1}{Online learning}{stream}{}{$\bullet$\ Why a validated model becomes obsolete\\[3pt]$\bullet$\ Offline versus online learning, and prequential evaluation\\[3pt]$\bullet$\ The streaming learning loop in River}

\begin{cframe}{Why a validated model becomes obsolete}{\cterm{Data-generating process}{the mechanism, physical or social, that produces the observations and their labels.}}
\centering
\begin{tikzpicture}
\infotile{0}{0}{6.6}{1.6}{tileblue}{courseblue}{users}{Behavior changes}{Customers, fraudsters, and operators adapt to new conditions and to the model itself.}
\infotile{6.8}{0}{6.6}{1.6}{tileblue}{courseblue}{thermometer-half}{Sensors degrade}{Probes lose calibration, are replaced, or are recalibrated with a different offset.}
\infotile{0}{-1.75}{6.6}{1.6}{tileblue}{courseblue}{gavel}{Rules change}{Business policies, regulations, and product formulations change the target itself.}
\infotile{6.8}{-1.75}{6.6}{1.6}{tileblue}{courseblue}{bug}{New errors appear}{Situations absent from the training data occur only after deployment.}
\keyidea{-3.65}{Adaptation does not indicate that the initial model was wrong: it acknowledges that the data-generating process evolves over time.}
\end{tikzpicture}
\coursefooter
\end{cframe}

\begin{cframe}{Offline versus online learning}{\cterm{Prequential evaluation}{test-then-train: each observation first tests the model and only then updates it.}}
\centering
\begin{tikzpicture}
\node[tileblue, anchor=north west, minimum width=6.6cm, minimum height=2.65cm] at (0,0) {};
\badge{0.6}{-0.55}{courseblue}{database}
\node[anchor=west, font=\small\bfseries, text=courseblue] at (1.15,-0.55) {Offline (batch)};
\node[anchor=north west, text width=6.1cm, align=left, font=\scriptsize, inner sep=0pt] at (0.25,-1.05) {$\bullet$\ Trained once on a fixed dataset $\mathcal{D} = \{(\mathbf{x}_i, y_i)\}_{i=1}^{n}$.\\[3pt]$\bullet$\ Benchmarked on a held-out test set.\\[3pt]$\bullet$\ Adapts only through periodic retraining.};
\node[tilegreen, anchor=north west, minimum width=6.6cm, minimum height=2.65cm] at (6.8,0) {};
\badge{7.4}{-0.55}{coursegreen}{stream}
\node[anchor=west, font=\small\bfseries, text=courseblue] at (7.95,-0.55) {Online (incremental)};
\node[anchor=north west, text width=6.1cm, align=left, font=\scriptsize, inner sep=0pt] at (7.05,-1.05) {$\bullet$\ Updated after each observation $(\mathbf{x}_t, y_t)$, with constant memory.\\[3pt]$\bullet$\ Operates directly on unbounded data streams.\\[3pt]$\bullet$\ Adapts continuously to change.};
\node[tilegold, anchor=north west, minimum width=13.4cm, minimum height=1.35cm, text width=12.8cm, align=left, font=\footnotesize] at (0,-2.85) {{\bfseries\goldink\faIcon{redo-alt}\ Prequential evaluation.} At time $t$ the model $f_{t-1}$ predicts $\hat{y}_t = f_{t-1}(\mathbf{x}_t)$, the loss $\ell(y_t, \hat{y}_t)$ is accumulated, and only then is $f_t$ obtained from $(\mathbf{x}_t, y_t)$. Every prediction is therefore made on an unseen observation, and no held-out set is required.};
\keyidea{-4.35}{The stream is both the training set and the test set; their order in time is what keeps the evaluation honest.}
\end{tikzpicture}
\coursefooter
\end{cframe}

\begin{cframe}{The streaming learning loop}{\cterm{River}{a Python library for online machine learning: every estimator learns one observation at a time.}\cterm{SGD}{stochastic gradient descent: a parameter update from the loss of a single observation.}}
\centering
\begin{tikzpicture}
\foreach \x/\ico/\ttl/\txt/\col in {0/satellite-dish/Observe/{$\mathbf{x}_t$ arrives}/courseblue, 2.55/eye/Predict/{\texttt{predict\_one}}/coursegold, 5.1/chart-line/Evaluate/{\texttt{metric.update}}/coursegold, 7.65/graduation-cap/Learn/{\texttt{learn\_one}}/coursegreen, 10.2/search/Check drift/{\texttt{detector.update}}/coursered} {
  \node[tile, fill=white, draw=\col, thick, anchor=north west, minimum width=2.3cm, minimum height=1.5cm] at (\x,0) {};
  \node[text=\col, font=\large] at (\x+1.15,-0.4) {\faIcon{\ico}};
  \node[font=\scriptsize\bfseries, text=courseblue] at (\x+1.15,-0.85) {\ttl};
  \node[font=\scriptsize] at (\x+1.15,-1.2) {\txt};
}
\foreach \x in {2.3,4.85,7.4,9.95} {\draw[arrowgold] (\x,-0.75) -- (\x+0.25,-0.75);}
\draw[arrowgray] (11.35,0) -- (11.35,0.25) -- (1.15,0.25) -- (1.15,0);
\node[font=\tiny\itshape, text=coursegray, fill=white, inner sep=1pt] at (6.25,0.25) {next observation, $y_t$ available after the prediction};
\node[anchor=west, font=\tiny\bfseries, text=coursegray] at (0,-1.85) {THE REPOSITORY EXAMPLE: \texttt{notebooks/Class8\_Model\_Incremental\_learning.ipynb}};
\node[tileblue, anchor=north west, minimum width=6.6cm, minimum height=1.95cm, text width=6.2cm, align=left, font=\scriptsize] at (0,-2.05) {{\bfseries\color{courseblue}\faIcon{cloud-sun}\ The stream.} The current temperature in Paris is predicted from the wind speed, both obtained from the Open-Meteo service. The service refreshes its values every 15 minutes, so faster polling only repeats the same observation.};
\node[tileblue, anchor=north west, minimum width=6.6cm, minimum height=1.95cm, text width=6.2cm, align=left, font=\scriptsize] at (6.8,-2.05) {{\bfseries\color{courseblue}\faIcon{project-diagram}\ The model.} A pipeline \texttt{StandardScaler | LinearRegression}, whose scaler statistics and weights (SGD, learning rate 0.01) are updated at every observation; the metric is the mean absolute error (MAE).};
\keyidea{-4.2}{\texttt{predict\_one} always precedes \texttt{learn\_one}: this ordering is prequential evaluation, written as code.}
\end{tikzpicture}
\coursefooter
\end{cframe}

% ================================================================== PART 2
\sectiondivider{2}{Concept drift}{random}{an online model learns from every observation, and is evaluated prequentially on the same stream.}{$\bullet$\ Definition: virtual and real drift\\[3pt]$\bullet$\ Four temporal patterns of drift\\[3pt]$\bullet$\ Drift in the two running cases}

\begin{cframe}{What is concept drift?}{\cterm{Covariate shift}{a change of the input distribution $P(\mathbf{x})$ only.}}
\centering
\begin{tikzpicture}
\node[tileblue, anchor=north west, minimum width=13.4cm, minimum height=1.35cm, text width=12.8cm, align=center] at (0,0) {{\bfseries\color{courseblue}Concept drift} is a change over time of the joint distribution:\quad $P_t(\mathbf{x}, y) \neq P_{t+\Delta}(\mathbf{x}, y)$, \quad with \quad $P(\mathbf{x}, y) = P(\mathbf{x})\,P(y \mid \mathbf{x})$.};
\infotile{0}{-1.55}{6.6}{2.1}{tilegold}{coursegold}{arrows-alt}{Virtual drift}{$P(\mathbf{x})$ changes, $P(y \mid \mathbf{x})$ does not (covariate shift). The boundary stays correct, but is used where training samples were scarce.}
\infotile{6.8}{-1.55}{6.6}{2.1}{tilered}{coursered}{exchange-alt}{Real drift}{$P(y \mid \mathbf{x})$ changes. The learned decision function becomes incorrect, whether or not the inputs change.}
\warning{-3.85}{exclamation-triangle}{Caution.}{A model trained before a real drift may still produce confident predictions that are systematically wrong. Drift is a change of the process, not noise in individual readings.}
\end{tikzpicture}
\coursefooter
\end{cframe}

\begin{cframe}{Virtual versus real drift}{\cterm{Error stream}{the sequence of prediction errors, available once the labels arrive.}}
\centering
\begin{tikzpicture}
\node[anchor=north west, inner sep=0pt] at (0.6,0.1) {\includegraphics[width=12.2cm]{images/class19_virtual_real_drift.pdf}};
\node[anchor=north west, font=\tiny\itshape, text=coursegray] at (0.6,-3.7) {illustrative synthetic data};
\node[tilegold, anchor=north west, minimum width=4.35cm, minimum height=1.25cm, text width=4.05cm, align=left, font=\scriptsize] at (0,-4.0) {{\bfseries\color{courseblue}Virtual.} The inputs move; the class of each point is unchanged, and so is the boundary.};
\node[tilered, anchor=north west, minimum width=4.35cm, minimum height=1.25cm, text width=4.05cm, align=left, font=\scriptsize] at (4.52,-4.0) {{\bfseries\color{courseblue}Real.} The same inputs now belong to the other class: the boundary itself is wrong.};
\node[tileblue, anchor=north west, minimum width=4.35cm, minimum height=1.25cm, text width=4.05cm, align=left, font=\scriptsize] at (9.05,-4.0) {{\bfseries\color{courseblue}Consequence.} Input monitoring detects virtual drift; real drift requires the labels.};
\end{tikzpicture}
\coursefooter
\end{cframe}

\begin{cframe}{Temporal patterns of drift}{\cterm{Recurring drift}{a previous concept returns, for example with the seasons.}}
\centering
\begin{tikzpicture}
\node[anchor=north west, inner sep=0pt] at (0,0.1) {\includegraphics[width=13.4cm]{images/class19_drift_types.pdf}};
\node[anchor=north west, font=\tiny\itshape, text=coursegray] at (0,-2.95) {illustrative synthetic data; dots are observations, lines the underlying concept};
\foreach \i/\col/\ttl/\txt in {0/coursered/Sudden/{The concept is replaced at once: a sensor recalibration, a change of product.}, 1/coursegold/Gradual/{Old and new concepts alternate; the new one becomes more frequent.}, 2/courseblue/Incremental/{The concept moves through intermediate states: wear of heating equipment.}, 3/coursegreen/Recurring/{A previous concept returns: seasonal patterns that can reuse earlier knowledge.}} {
  \pgfmathsetmacro{\x}{\i*3.4}
  \node[tile, fill=white, draw=\col, thick, anchor=north west, minimum width=3.25cm, minimum height=1.55cm, text width=3.0cm, align=left, font=\scriptsize] at (\x,-3.3) {{\bfseries\color{courseblue}\ttl}\\[1pt]\txt};
}
\end{tikzpicture}
\coursefooter
\end{cframe}

% ================================================================== PART 3
\sectiondivider{3}{Detection}{search}{drift is a change of $P(\mathbf{x})$, of $P(y \mid \mathbf{x})$, or of both, and it unfolds suddenly, gradually, incrementally, or recurrently.}{$\bullet$\ How ADWIN decides that a mean has changed\\[3pt]$\bullet$\ ADWIN on the pasteurization stream\\[3pt]$\bullet$\ Detectors and drift reports together}

\begin{cframe}{How ADWIN works}{\cterm{ADWIN}{adaptive windowing (Bifet and Gavald\`a, 2007): a change detector with a window of variable length.}\cterm{Hoeffding bound}{a bound on how far a sample mean can lie from its expectation, with a given confidence.}}
\centering
\begin{tikzpicture}
\node[anchor=west, font=\tiny\bfseries, text=coursegray] at (0,0.15) {THE WINDOW $W$ OF RECENT VALUES (HERE, ABSOLUTE ERRORS), SPLIT AT EVERY POSITION};
\foreach \k in {0,...,21} {
  \pgfmathsetmacro{\x}{\k*0.42}
  \ifnum\k<13 \def\wcol{courseblue!55}\else \def\wcol{coursered!70}\fi
  \node[rounded corners=1pt, fill=\wcol, minimum width=0.36cm, minimum height=0.5cm, inner sep=0pt, anchor=north west] at (\x,-0.15) {};
}
\draw[coursegold, line width=1.6pt] (5.44,0.0) -- (5.44,-0.85);
\draw[decorate, decoration={brace, amplitude=4pt, mirror}, thick, courseblue] (0,-0.75) -- (5.38,-0.75) node[midway, below=5pt, font=\scriptsize] {$W_0$: older, $n_0$ values, mean $\hat{\mu}_{W_0}$};
\draw[decorate, decoration={brace, amplitude=4pt, mirror}, thick, coursered] (5.5,-0.75) -- (9.18,-0.75) node[midway, below=5pt, font=\scriptsize] {$W_1$: recent, $n_1$ values, $\hat{\mu}_{W_1}$};
\node[tileblue, anchor=north west, minimum width=3.95cm, minimum height=1.45cm, text width=3.7cm, align=left, font=\scriptsize] at (9.45,0.05) {{\bfseries\color{courseblue}Decision.} If the two means differ by at least $\varepsilon_{\mathrm{cut}}$, a change is detected and $W_0$ is dropped.};
\node[tilegold, anchor=north west, minimum width=13.4cm, minimum height=1.25cm, text width=12.8cm, align=center] at (0,-1.6) {$|\hat{\mu}_{W_0} - \hat{\mu}_{W_1}| \ge \varepsilon_{\mathrm{cut}}, \qquad \varepsilon_{\mathrm{cut}} = \sqrt{\frac{1}{2m}\ln\frac{4}{\delta'}}, \qquad m = \frac{1}{1/n_0 + 1/n_1}, \qquad \delta' = \frac{\delta}{n}$};
\infotile{0}{-3.0}{6.6}{1.6}{tileblue}{courseblue}{sliders-h}{The confidence parameter $\delta$}{It controls the false-alarm rate; here $\delta = 0.002$.}
\infotile{6.8}{-3.0}{6.6}{1.6}{tilegreen}{coursegreen}{expand-alt}{An adaptive window}{It grows while the stream is stable; it shrinks at a change.}
\keyidea{-4.75}{A statistically grounded trigger for retraining or resetting, instead of a fixed schedule.}
\end{tikzpicture}
\coursefooter
\end{cframe}

\begin{cframe}{ADWIN on the pasteurization stream}{\cterm{Soft sensor}{a model that estimates a quantity from other, more easily measured quantities.}\cterm{Detection delay}{the number of samples between the occurrence of a drift and its detection.}}
\centering
\begin{tikzpicture}
\node[tileblue, anchor=north west, minimum width=13.4cm, minimum height=1.05cm, text width=12.8cm, align=left, font=\scriptsize] at (0,0.1) {{\bfseries\color{courseblue}\faIcon{flask}\ Experiment.} An online soft sensor estimates the viscosity $\eta$ from the temperature $T$ (60 batches, 19{,}530 samples at 1\,Hz). After batch 30 a product change raises $\eta$ by 0.4\,cP at every temperature: a sudden real drift.};
\node[anchor=north west, inner sep=0pt] at (0,-1.15) {\includegraphics[width=9.0cm]{images/class19_adwin_detection.pdf}};
\node[tilegreen, anchor=north west, minimum width=4.2cm, minimum height=1.45cm, text width=3.95cm, align=left, font=\scriptsize] at (9.2,-1.15) {{\bfseries\color{courseblue}Delay: 51 samples.}\\[1pt]The window grows while the error stream is stationary and collapses 51\,s after the drift.};
\node[tileblue, anchor=north west, minimum width=4.2cm, minimum height=1.45cm, text width=3.95cm, align=left, font=\scriptsize] at (9.2,-2.75) {{\bfseries\color{courseblue}One early alarm.}\\[1pt]Apart from the initial learning phase, no false alarm in 19{,}530 samples.};
\node[anchor=north west, font=\tiny\itshape, text=coursegray, text width=13.4cm] at (0,-4.3) {Data from \texttt{pasteurization/synth\_sensors.py} with an injected drift; River 0.22, \texttt{LinearRegression} with SGD, \texttt{ADWIN(delta=0.002)}.};
\end{tikzpicture}
\coursefooter
\end{cframe}

\begin{cframe}{Monitoring beyond a single detector}{\cterm{Drift report}{a periodic statistical comparison of every input feature against a reference distribution.}}
\centering
\begin{tikzpicture}
\infotile{0}{0}{6.6}{2.5}{tilered}{coursered}{search}{ADWIN on the error stream}{Fast and targeted: it detects a degradation of performance as soon as the labels arrive. It sees real drift and the virtual drift that hurts the model.}
\infotile{6.8}{0}{6.6}{2.5}{tileblue}{courseblue}{file-medical-alt}{Drift report on the inputs}{Comprehensive: it tests the distribution of every feature (Evidently \texttt{DataDriftPreset}, \texttt{DataQualityPreset}) and can reveal virtual drift before performance degrades.}
\node[tile, fill=courseice, anchor=north west, minimum width=13.4cm, minimum height=0.95cm, text width=13.0cm, align=left, font=\scriptsize] at (0,-2.7) {{\bfseries\color{courseblue}\faIcon{book}\ In the repository.} \texttt{notebooks/Class8\_Model\_Incremental\_learning.ipynb} produces both: River with ADWIN on the error stream, and Evidently reports on the inputs.};
\keyidea{-3.85}{A common production pattern combines a lightweight online detector with periodic, more comprehensive reports.}
\end{tikzpicture}
\coursefooter
\end{cframe}

% ================================================================== PART 4
\sectiondivider{4}{Adaptation}{sync-alt}{ADWIN detects a change of the mean of the error stream with a controlled false-alarm rate.}{$\bullet$\ Three adaptation strategies, compared on the same stream\\[3pt]$\bullet$\ Batch retraining versus online learning\\[3pt]$\bullet$\ Where to find all of this in the repository}

\begin{cframe}{Adaptation strategies compared}{\cterm{Reset loop}{a sequence of resets in which each untrained model triggers the next detection.}\cterm{Warm start}{continuing to train the current model instead of starting from scratch.}}
\centering
\begin{tikzpicture}
\node[anchor=north west, inner sep=0pt] at (0,0.1) {\includegraphics[width=8.4cm]{images/class19_adaptation_strategies.pdf}};
\node[anchor=west, font=\tiny\bfseries, text=coursegray] at (8.65,0.0) {MAE BEFORE $\to$ AFTER THE DRIFT};
\foreach \y/\col/\ttl/\txt in {-0.25/coursegray/{Batch model, frozen}/{0.041 $\to$ 0.398: no adaptation.}, -1.25/coursegreen/{Online, no reset}/{0.052 $\to$ 0.043: recovers in 200 samples.}, -2.25/coursered/{Online, reset at detection}/{0.599 $\to$ 0.820: 159 resets.}} {
  \node[tile, fill=white, draw=\col, thick, anchor=north west, minimum width=4.75cm, minimum height=0.92cm, text width=4.5cm, align=left, font=\scriptsize, inner sep=3pt] at (8.65,\y) {{\bfseries\color{courseblue}\ttl}\\\txt};
}
\warning{-3.4}{redo}{Why the reset fails.}{An untrained model has a large and decreasing error, which is itself a change that ADWIN detects: the model never stabilizes. Safer responses are to warm-start from the current model, retrain on the recent window $W_1$, or only raise an alert.}
\end{tikzpicture}
\coursefooter
\end{cframe}

\begin{cframe}{Batch retraining versus online learning}{\cterm{Auditability}{the possibility of reviewing, afterwards, which model made a decision and why.}}
\centering
\begin{tikzpicture}
\node[chip, anchor=north west, minimum width=5.05cm, minimum height=0.42cm, font=\tiny\bfseries] at (3.15,0) {BATCH RETRAINING};
\node[chip, fill=coursegreen, anchor=north west, minimum width=5.05cm, minimum height=0.42cm, font=\tiny\bfseries] at (8.35,0) {ONLINE LEARNING};
\comparerow{-0.55}{stopwatch}{Adaptation}{Waits for the next retraining cycle.}{Adapts at every observation.}
\comparerow{-1.35}{clipboard-check}{Auditability}{Fixed snapshots that can be reviewed and versioned.}{The model state changes continuously.}
\comparerow{-2.15}{balance-scale}{Stability}{Stable between cycles.}{May overreact to noise if the detector or learning rate is too sensitive.}
\comparerow{-2.95}{server}{Infrastructure}{Large periodic training jobs.}{An always-on service, with small updates.}
\keyidea{-3.95}{The choice is a trade-off between the speed of adaptation and the control over what is deployed; many systems combine both.}
\end{tikzpicture}
\coursefooter
\end{cframe}

\begin{cframe}{Where to find it in the repository}{\cterm{Scenario}{a reproducible fault injected by the classroom simulator.}}
\centering
\begin{tikzpicture}
\node[anchor=west, font=\tiny\bfseries, text=coursegray] at (0,0.15) {WHERE TO LOOK};
\node[anchor=west, font=\tiny\bfseries, text=coursegold!65!black] at (4.75,0.15) {WHAT IT SHOWS};
\node[anchor=west, font=\tiny\bfseries, text=coursegray] at (11.25,0.15) {USED IN};
\reporow{-0.3}{notebooks/Class8\_Model\_\\Incremental\_learning.ipynb}{River, ADWIN, and Evidently on a weather stream.}{courseblue}{Class 29}
\reporow{-1.12}{streamlit/pages/App\_Stream.py}{The same pattern as a live dashboard.}{courseblue}{Class 29}
\reporow{-1.94}{pasteurization/serving.py}{The simulator; \texttt{sensor\_drift} injects a calibration drift.}{coursegreen}{Class 30}
\reporow{-2.76}{presentations/generate\_\\class30\_figures.py}{An online baseline and ADWIN on five fault scenarios.}{coursegreen}{Class 30}
\keyidea{-3.8}{Class 30 turns this class into a feature of the challenge application: learn normal behavior online and record drift events.}
\end{tikzpicture}
\coursefooter
\end{cframe}

% ------------------------------------------------------------------ closing
\section{Wrap-up}\gdef\currentsectionname{Wrap-up}
\begin{frame}{From this class to the next}
\centering
\begin{tikzpicture}
\foreach \i/\lab/\bg/\fg in {1/Problem/courseice/coursegray, 2/Data/courseice/coursegray, 3/Features/courseice/coursegray, 4/Model/courseblue/white, 5/Interface/courseice/coursegray, 6/Delivery/courseice/coursegray, 7/Monitor/courseblue/white, 8/Evolution/courseblue/white} {
  \node[rounded corners=3pt, minimum width=1.55cm, minimum height=0.6cm, font=\scriptsize\bfseries, fill=\bg, text=\fg] at ({0.78+(\i-1)*1.68},0) {\lab};
}
\node[font=\tiny\bfseries, text=courseblue] at (6.7,-0.52) {THIS CLASS CONNECTS THE MODEL TO MONITORING AND EVOLUTION};
\foreach \x/\sty/\ico/\ttl/\txt in {0/tileblue/lightbulb/What stays with us/{$\bullet$\ Evaluate prequentially.\\[3pt]$\bullet$\ Detect drift with a statistical test.\\[3pt]$\bullet$\ Define in advance how the model is updated after detection.}, 4.55/tileblue/project-diagram/Use it in your project/{$\bullet$\ Say which drift your system can observe: virtual, real, or both.\\[3pt]$\bullet$\ Report detection delay and false alarms, not only accuracy.}, 9.1/tilegold/arrow-right/Next: Class 30/{\textbf{Pasteurization Challenge, Session 3.}\\[3pt]A second change request: the application learns normal behavior online and records drift events.}} {
  \node[\sty, anchor=north west, minimum width=4.3cm, minimum height=3.8cm] at (\x,-0.95) {};
  \badge{\x+0.65}{-1.6}{courseblue}{\ico}
  \node[anchor=west, text width=3.0cm, align=left, font=\small\bfseries, text=courseblue] at (\x+1.2,-1.6) {\ttl};
  \node[anchor=north west, text width=3.9cm, align=left, font=\scriptsize] at (\x+0.2,-2.25) {\txt};
}
\node[tilegold, fill=white, anchor=north west, minimum width=13.4cm, minimum height=0.95cm, text width=12.6cm, align=left, font=\footnotesize] at (0,-5.0) {{\bfseries\goldink\faIcon{route}\ Course thread.} The data-generating process changes over time; a deployed model needs a statistical way of noticing it.};
\end{tikzpicture}
\coursefooter
\end{frame}

% ------------------------------------------------------------------ references
\begin{frame}{Papers and materials}
\centering
\begin{tikzpicture}
\refrow{0}{2.35}{random}{Concept drift}{Gama, J., \v{Z}liobait\.{e}, I., Bifet, A., Pechenizkiy, M., and Bouchachia, A. (2014). A survey on concept drift adaptation. \emph{ACM Computing Surveys}, 46(4), 44:1--44:37.\\Lu, J., Liu, A., Dong, F., Gu, F., Gama, J., and Zhang, G. (2019). Learning under concept drift: a review. \emph{IEEE Transactions on Knowledge and Data Engineering}, 31(12), 2346--2363.}
\refrow{-2.35}{2.05}{search}{Detection}{Bifet, A. and Gavald\`a, R. (2007). Learning from time-changing data with adaptive windowing. \emph{Proceedings of the SIAM International Conference on Data Mining}, 443--448.\\Gama, J., Sebasti\~ao, R., and Rodrigues, P. P. (2013). On evaluating stream learning algorithms. \emph{Machine Learning}, 90(3), 317--346.}
\refrow{-4.4}{1.15}{stream}{Software}{Montiel, J. et al. (2021). River: machine learning for streaming data in Python. \emph{Journal of Machine Learning Research}, 22(110), 1--8.\\Documentation: \href{https://riverml.xyz/latest/}{River}; \href{https://docs.evidentlyai.com/}{Evidently}.}
\end{tikzpicture}
\coursefooter
\end{frame}

\end{document}
